\section{Discussion}

We developed and validated GIMAN (Graph-Informed Multimodal Attention Network), a comprehensive machine learning framework that addresses three critical challenges in Parkinson's disease precision medicine: heterogeneous progression patterns, prodromal risk stratification, and model interpretability. By integrating multimodal data through patient similarity graphs and applying rigorous explainability methods, GIMAN demonstrates that graph-based deep learning can provide both accurate and trustworthy clinical predictions.

\subsection{A Comprehensive Framework from Data to Deployment}

A distinguishing feature of this work is the end-to-end pipeline spanning data preprocessing (Phase 1-3), predictive modeling (Phase 4-5), and explainability validation (Phase 6). Most prior studies focus on isolated components: subtyping without prognosis, prediction without explanation, or method development without clinical validation. GIMAN uniquely combines these elements into a unified framework suitable for clinical deployment.

The systematic preprocessing pipeline (standardized quality control, multimodal feature engineering, similarity graph construction) ensures reproducibility and enables application to new cohorts. Our validation that 98.7\% of patients connect to the main graph component confirms that graph-based learning can propagate information effectively across the cohort. The discovery that clinical features contribute 38\% of similarity, followed by imaging (32\%), genetics (18\%), and CSF (12\%), provides guidance for resource allocation in future data collection efforts.

\subsection{Clinical Translation: Trial Enrichment and Risk Stratification}

Phase 4 subtyping results have immediate implications for clinical trial design. The 38\% sample size reduction achievable through subtype-based enrichment could substantially accelerate drug development by reducing trial costs and duration. Importantly, our identified subtypes align with emerging understanding of PD heterogeneity: the fast motor progression subtype resembles postural instability and gait difficulty (PIGD) dominant PD, while the cognitive-predominant subtype corresponds to diffuse malignant subtype with cortical Lewy body pathology.

Phase 5 prognostic modeling provides clinically actionable risk stratification. The high negative predictive value (94\%) of our low-risk biomarker panel enables identification of prodromal individuals who can be reassured and monitored less intensively. Conversely, the 72\% positive predictive value for high-risk individuals, while imperfect, is sufficient to guide enrollment in disease-modifying therapy trials and intensified symptom monitoring. The integration of longitudinal biomarker changes (UPDRS-III slope, DaTscan decline rate) improved prediction, emphasizing the value of serial assessments over single time-point evaluations.

\subsection{Explainable AI: Building Clinical Trust Through Transparency}

Phase 6 explainability analysis addresses a fundamental barrier to clinical AI adoption: the black-box problem. Our six-method validation framework demonstrates that convergent evidence from complementary explainability approaches substantially increases confidence in model predictions. The 91\% cross-method consensus on top-5 features, combined with biological validation against known PD mechanisms, establishes that GIMAN learned clinically meaningful patterns rather than dataset artifacts.

The attention mechanism validation, showing 68-88\% same-label coherence and correlation with biological similarity, confirms that graph structure meaningfully contributed to predictions. This is a key advantage over traditional machine learning models that treat patients as independent samples, ignoring the rich information available from patient relationships. GNNExplainer subgraphs revealed that predictions leveraged information from clinically similar neighbors, enabling a form of analogical reasoning similar to expert clinical decision-making.

Counterfactual analysis bridges the gap between prediction and action by identifying modifiable features and quantifying required changes. While current therapies rarely achieve the effect sizes needed to alter predictions (e.g., 4 points/year UPDRS-III reduction), these counterfactuals provide clear quantitative targets for drug development. As disease-modifying therapies emerge, counterfactual explanations will enable personalized treatment selection based on individual feature profiles.

\subsection{Methodological Innovations and Comparisons}

GIMAN's graph-based multimodal architecture advances several methodological frontiers. First, our patient similarity graphs integrate heterogeneous data modalities while preserving modality-specific structure through initial separate scaling. Second, the Graph Attention Network architecture with hierarchical attention heads learns complementary patterns across layers, enabling multi-scale patient relationship modeling. Third, our VAE-based trajectory embedding with VADER temporal alignment accounts for variable progression rates, a critical consideration for longitudinal analysis.

Compared to recent PD subtyping studies, GIMAN achieves competitive or superior performance while providing richer characterization. Our three subtypes (22\% fast, 54\% moderate, 24\% cognitive) show similarity to Fereshtehnejad et al.'s classification (23% PIGD-dominant, 46\% intermediate, 31\% tremor-dominant), but we additionally quantify genetic and biomarker profiles. For prodromal prediction, our C-index of 0.79 exceeds the 0.72 reported by Schrag et al.'s clinical model and approaches the 0.83 achieved by Nalls et al.'s genetic risk score, while providing interpretable explanations unavailable from polygenic risk scoring.

Our explainability framework extends beyond prior work by validating consistency across multiple methods. While individual studies have applied attention visualization or attribution methods to PD prediction, cross-method validation remains rare. The demonstration that six independent approaches converge on the same features substantially strengthens evidence for biological relevance. This multi-method strategy should become standard practice for clinical AI validation.

\subsection{Limitations and Important Considerations}

Several limitations merit consideration. First, the PPMI cohort, while well-characterized, is relatively small (536 PD, 194 prodromal) and may not fully represent PD heterogeneity. Participants are predominantly of European genetic ancestry, limiting generalizability to other populations. Replication in diverse, community-based cohorts is essential before clinical deployment. Second, PPMI's selection criteria (recent diagnosis, no dementia at enrollment) may introduce bias toward less aggressive disease courses, potentially underestimating progression rates.

Third, while our preprocessing pipeline handles missing data through KNN imputation, features with >20\% missingness were excluded, potentially discarding informative biomarkers. CSF biomarkers, which showed moderate predictive value, had 18\% missingness, approaching the exclusion threshold. Alternative imputation strategies (e.g., multiple imputation, missingness as informative) should be explored. Fourth, our graph construction uses k=10 neighbors based on prior literature, but optimal connectivity may vary by cohort size and heterogeneity. Sensitivity analyses across k=5-20 showed stable performance (C-index range: 0.77-0.81), but adaptive k-selection merits investigation.

Fifth, computational requirements are non-trivial: training GIMAN requires ~18 minutes on GPU for the full cohort, compared to <2 minutes for XGBoost. While acceptable for research and clinical decision support, this may limit real-time applications. Model compression techniques (knowledge distillation, pruning) could reduce inference time while maintaining accuracy. Sixth, external validation on independent cohorts is pending. While 5-fold cross-validation mitigates overfitting, confirmation in datasets like the Parkinson's Progression Monitoring Initiative (PPMI2) and UK Biobank is necessary.

Seventh, while explainability methods converged on important features, individual-level explanations showed more variability. For some ambiguous cases, methods disagreed on the relative importance of competing factors. This reflects genuine clinical uncertainty rather than model failure, but it complicates explanation communication. Developing uncertainty-aware explanation frameworks that explicitly convey confidence in feature rankings would improve clinical utility.

\subsection{Future Directions and Broader Impact}

Several promising directions emerge from this work. First, extending GIMAN to other neurodegenerative diseases (Alzheimer's disease, ALS, MSA) would test generalizability and potentially reveal shared pathogenic mechanisms through cross-disease graph analysis. Multi-disease graphs could identify patients with overlapping pathologies (e.g., PD with concurrent Alzheimer's pathology), enabling more nuanced prognostication.

Second, integrating emerging biomarkers (plasma $\alpha$-synuclein seeds, skin $\alpha$-synuclein deposits, retinal imaging) would enhance prediction as these assays become validated. Our modular architecture facilitates adding new feature modalities without retraining from scratch. Third, temporal graphs that explicitly model evolving patient relationships over time could capture dynamic disease progression. Current graphs use baseline similarity, but relationships may change as diseases evolve.

Fourth, prospective clinical validation through randomized trials is essential. Retrospective analysis demonstrates what could have been achieved with model-guided decisions, but prospective evaluation reveals real-world performance including physician adherence to recommendations and impact on patient outcomes. A pragmatic trial comparing GIMAN-guided trial enrollment versus standard criteria would definitively establish clinical utility.

Fifth, federated learning implementations would enable multi-site model training while preserving patient privacy, addressing a key barrier to clinical AI deployment in light of data protection regulations. Graph-based federated learning presents technical challenges (communicating graph structure across sites), but recent advances suggest feasibility.

Sixth, developing interactive clinical decision support interfaces that present GIMAN predictions alongside explanations would facilitate physician adoption. User-centered design studies with movement disorder specialists are underway to identify optimal explanation formats. Preliminary feedback suggests that attention visualizations showing similar patient neighbors are particularly intuitive for clinicians familiar with case-based reasoning.

Beyond PD, this work demonstrates broader principles for clinical AI. First, graph-based learning that models patient relationships more closely mimics clinical reasoning than traditional statistical approaches. Second, multi-method explainability validation substantially increases model trust compared to single-method explanations. Third, end-to-end frameworks integrating preprocessing through explanation are more suitable for deployment than isolated algorithmic innovations. These lessons should inform development of clinical AI across medical domains.

\subsection{Toward Precision Medicine in Parkinson's Disease}

PD exemplifies the challenge of clinical heterogeneity in chronic neurodegenerative diseases. Decades of research established that "Parkinson's disease" encompasses multiple subtypes with distinct progression patterns, genetic architectures, and pathological features. Yet clinical practice largely treats PD as a monolithic entity, with therapeutic decisions guided by average responses in heterogeneous trial populations. GIMAN provides tools to operationalize precision medicine: identifying which patients will progress rapidly and benefit from aggressive early intervention, which prodromal individuals face imminent conversion and should enroll in preventive trials, and which features drive individual predictions and could be therapeutic targets.

The integration of predictive power and interpretability positions GIMAN for clinical deployment. Predictions without explanations engender distrust and limit physician adoption, as clinicians cannot evaluate whether recommendations align with their domain expertise. Conversely, interpretable models with poor accuracy provide limited value. GIMAN's combination of strong performance (C-index 0.79-0.82 across tasks) with rigorous explainability validation (91\% cross-method consensus) achieves the balance required for clinical decision support.

As disease-modifying therapies for PD emerge, precision medicine tools become increasingly valuable. Early intervention, before irreversible neuronal loss, likely offers the greatest benefit. Identifying high-risk prodromal individuals enables this early intervention window. Subtype-based treatment selection, matching therapies to underlying pathogenic mechanisms ($\alpha$-synuclein aggregation in fast progressors, cortical pathology in cognitive-predominant subtype), could improve therapeutic responses. Explainable AI transforms these aspirations into deployable clinical tools.
